%% ============================================================
%%  CSCI 4922/5922 — Lab Assignment 3 Report
%%  Neural Networks and Deep Learning, Spring 2026
%%  Visual Question Answering on VizWiz
%% ============================================================

\documentclass[11pt]{article}

\usepackage[margin=1in]{geometry}
\usepackage{amsmath, amssymb}
\usepackage{booktabs}
\usepackage{graphicx}
\usepackage{hyperref}
\usepackage{caption}
\usepackage{subcaption}
\usepackage{enumitem}
\usepackage{titlesec}
\usepackage{parskip}
\usepackage{listings}
\usepackage{xcolor}
\usepackage{microtype}

\lstset{
  language=Python,
  basicstyle=\ttfamily\footnotesize,
  keywordstyle=\color{blue},
  commentstyle=\color{gray},
  stringstyle=\color{orange!80!black},
  breaklines=true,
  frame=single,
  numbers=left,
  numberstyle=\tiny\color{gray},
  xleftmargin=1.5em,
}

\title{
  \textbf{Lab Assignment 3: Visual Question Answering on VizWiz}\\
  \large CSCI 5922 — Neural Networks and Deep Learning, Spring 2026
}
\author{Rajvardhan Patil}
\date{}

\begin{document}
\maketitle

% ════════════════════════════════════════════════════════════
\section{Custom Multi-Modal Architectures (Challenges 1 \& 2)}
% ════════════════════════════════════════════════════════════

\subsection{Methods}

Images are resized to $112\times112$, randomly flipped horizontally during training, and
normalized with ImageNet statistics. Questions are lowercased, tokenized, and padded to
20 tokens with a 5{,}002-word vocabulary. For Task 1, class-weighted cross-entropy loss
counters the 73\% answerable imbalance ($w_0 \approx 2.71$, $w_1 = 1.0$). For Task 2,
the majority-vote answer over 10 responses is the training target, using a closed
vocabulary of the top-1{,}000 answers. Training uses the first 10{,}000 samples with images
pre-cached in RAM to avoid repeated disk I/O.

\paragraph{Experiment 3.1 — Attention Heads Ablation (Task 1).}
\textit{Hypothesis:} Performance will peak at 4 heads. With $d=256$, an 8-head config gives
only 32 dims per head, which may be too narrow; 4 heads at 64 dims each should be optimal.
\textit{Variable:} nhead $\in \{1, 2, 4, 8\}$ in the cross-modal fusion layer.
All other settings fixed: AdamW, lr $= 3\times10^{-4}$, batch 128, 15 epochs.

\paragraph{Experiment 3.2 — Supervision Strategy (Task 2).}
\textit{Hypothesis:} Majority-vote supervision will outperform using all 10 answers.
Training on all 10 answers per sample expands the dataset 10$\times$ but introduces
conflicting gradients when annotators disagree.
\textit{Variable:} Condition A — majority vote; Condition B — all 10 answers as separate examples.

\paragraph{Architecture.}
\textbf{Image:} 4-block CNN (Conv+BN+ReLU+Pool) trained from scratch, $112\to1\times1$,
projected to 256-d. \textbf{Text:} Embedding + sinusoidal PE + 2-layer Transformer,
mean-pooled to 256-d. \textbf{Fusion:} Image and text stacked as a 2-token sequence,
processed by multi-head self-attention, concatenated and projected to 256-d.
\textbf{Head (Task 1):} Linear(256,128) $\to$ ReLU $\to$ Dropout(0.3) $\to$ Linear(128,2).
\textbf{Head (Task 2):} Same but output is Linear(128,1000) over the answer vocabulary.

\subsection{Results}

\begin{table}[h]
  \centering
  \caption{Exp.\ 3.1: Fusion attention heads vs.\ validation accuracy.}
  \label{tab:exp31_cls}
  \begin{tabular}{lcc}
    \toprule
    \textbf{Condition} & \textbf{Val Acc (\%)} & \textbf{Notes} \\
    \midrule
    nhead = 1 & \textbf{74.11} & Widest per-head dim \\
    nhead = 2 & 73.74 & Marginal decrease \\
    nhead = 4 & 73.51 & Slight drop \\
    nhead = 8 & 74.32 & Recovers; statistically equivalent to nhead=1 \\
    \bottomrule
  \end{tabular}
\end{table}

\begin{table}[h]
  \centering
  \caption{Exp.\ 3.2: Supervision strategy vs.\ human answer accuracy.}
  \label{tab:exp32_gen}
  \begin{tabular}{lcc}
    \toprule
    \textbf{Condition} & \textbf{Val Answer Acc} & \textbf{Notes} \\
    \midrule
    A: Majority Vote  & \textbf{0.21} & Stable, consistent supervision \\
    B: All 10 Answers & 0.18          & Contradictory gradients hurt learning \\
    \bottomrule
  \end{tabular}
\end{table}

\subsection{Analysis}

\textbf{Task 1.} The hypothesis was not supported — all four head counts land within 0.8
points of each other. The fusion module operates on a 2-token sequence, so attention is
just a $2\times2$ matrix; multiple heads provide no complementary specialization here. We
select nhead=4 for the final model. The $\approx$74\% accuracy exceeds the 68\%
always-answerable baseline, confirming real signal is learned.

\textbf{Task 2.} Majority vote outperforms all-answers augmentation as predicted. When
one annotator says ``red'' and nine say ``unanswerable,'' training on all 10 creates
conflicting gradients. Majority vote avoids this. Overall answer accuracy is low (0.21)
because the VizWiz images are often blurry and the CNN is trained from scratch.

% ════════════════════════════════════════════════════════════
\section{Foundation Model Features: CLIP (Challenges 3 \& 4)}
% ════════════════════════════════════════════════════════════

\subsection{Methods}

Pre-extracted 512-d CLIP features are loaded directly for train, val, and test splits.
No additional normalization is needed since CLIP embeddings are already unit-normalized.
The 10{,}000 sample cap is enforced by slicing the training tensor.

\paragraph{Experiment 4.1 — Fusion Strategy (Task 1).}
\textit{Hypothesis:} Concatenation will outperform sum and cross-attention. Concat gives
the MLP independent access to both feature sets; sum discards channel identities;
cross-attention on length-1 sequences degenerates to a scalar weighting.
\textit{Variable:} Concat $[f_\text{img} \| f_\text{txt}]$ vs.\ Sum $f_\text{img}+f_\text{txt}$
vs.\ Cross-Attention. Fixed: Adam, lr $=3\times10^{-4}$, batch 256, 30 epochs.

\paragraph{Experiment 4.2 — Vocab Size (Task 2).}
\textit{Hypothesis:} Vocab=1{,}000 will perform best. Vocab=500 is too coarse; vocab=2{,}000
gives only $\approx$5 training examples per class, too few for reliable learning.
\textit{Variable:} Answer vocabulary size $\in \{500, 1000, 2000\}$.

\paragraph{Architecture.}
\textbf{Task 1 (Challenge 3):} Concat $\to$ Linear(1024,512)+LN+ReLU+Drop(0.3) $\to$
Linear(512,256)+LN+ReLU+Drop(0.21) $\to$ Linear(256,2), class-weighted loss.
\textbf{Task 2 (Challenge 4):} Same hidden layers, output Linear(256,1000), standard CE loss.

\subsection{Results}

\begin{table}[h]
  \centering
  \caption{Exp.\ 4.1: CLIP fusion strategy vs.\ binary classification accuracy.}
  \label{tab:exp41_cls}
  \begin{tabular}{lcc}
    \toprule
    \textbf{Condition} & \textbf{Val Acc (\%)} & \textbf{Notes} \\
    \midrule
    A: Concat $\to$ MLP    & \textbf{81.52} & Best; full access to both feature sets \\
    B: Sum $\to$ MLP       & 80.53 & Competitive; CLIP alignment helps \\
    C: Cross-Attention     & 69.44 & Degenerate for length-1 sequences \\
    \bottomrule
  \end{tabular}
\end{table}

\begin{table}[h]
  \centering
  \caption{Exp.\ 4.2: Answer vocabulary size vs.\ validation classification accuracy.}
  \label{tab:exp42_gen}
  \begin{tabular}{lcc}
    \toprule
    \textbf{Condition} & \textbf{Val Cls Acc} & \textbf{Notes} \\
    \midrule
    Vocab = 500  & \textbf{68.88} & Easier task; fewer classes \\
    Vocab = 1000 & 65.27 & Harder; data spread thinner \\
    Vocab = 2000 & 62.12 & Too sparse; $\approx$5 examples/class \\
    \bottomrule
  \end{tabular}
\end{table}

\subsection{Analysis}

\textbf{Task 1.} Concatenation wins as predicted. The MLP can learn image-only and
text-only signals independently before combining them. Sum conflates the two modalities.
Cross-attention on single-vector inputs is degenerate since attention collapses to a
scalar weight, making it equivalent to a linear combination.

\textbf{Task 2.} The 500-class vocab achieves the highest accuracy, contradicting the
hypothesis. As vocab grows, the classification task gets harder and per-class data
density drops too low. That said, classification accuracy is a proxy — a larger vocab
may cover more answer types even if harder to optimize.

\textbf{E2E vs.\ CLIP.} CLIP (concat) reaches 81.52\% vs.\ $\approx$74\% for E2E on
binary classification, a gap of 7--8 points. CLIP was pretrained on 400 million
image-text pairs and arrives with rich aligned representations. Our E2E model must
learn features, fusion, and classification from 10{,}000 samples from scratch, which is
a much harder problem. CLIP is the clear choice for low-data settings.

% ════════════════════════════════════════════════════════════
\section{Data Efficiency (CSCI 5922 Only)}
% ════════════════════════════════════════════════════════════

\subsection{Hypothesis}

The E2E model will benefit more from additional data than the CLIP model. The E2E CNN
must learn visual features from scratch and is likely data-limited at 50\%. The CLIP
MLP starts from semantically rich embeddings and should saturate earlier, since the
bottleneck is the embedding quality rather than labeled data quantity.

\subsection{Methods}

Both models use fixed hyperparameters from their competition runs. Training fractions:
50\% (10{,}000 samples), 75\% (15{,}392 samples), 100\% (20{,}523 samples). All runs use
seed 42; evaluation on the full 4{,}319-sample validation set.

\subsection{Results}

\begin{table}[h]
  \centering
  \caption{Validation accuracy across training data fractions.}
  \label{tab:data_eff}
  \begin{tabular}{llccc}
    \toprule
    \textbf{Model} & \textbf{Architecture} & \textbf{50\%} & \textbf{75\%} & \textbf{100\%} \\
    \midrule
    End-to-End & CNN + Transformer + Fusion & 74.3 & 73.7 & 75.4 \\
    CLIP-based & MLP on CLIP features       & 81.4 & 82.0 & 83.3 \\
    \bottomrule
  \end{tabular}
\end{table}

\subsection{Analysis}

The E2E model is non-monotonic: accuracy dips at 75\% before recovering at 100\%. This
reflects training variance in a from-scratch CNN — the 15K regime appears to be an
unstable midpoint where the model has seen more variety but not yet enough signal to
exploit it. Net gain from 50\% to 100\%: $+1.1$ points.

The CLIP model improves monotonically ($+1.9$ points), larger in absolute terms than the
E2E gain — partially contradicting the hypothesis. The CLIP MLP continues to sharpen
its decision boundary with more data rather than saturating.

The core hypothesis still holds in a broader sense: the CLIP model at 50\% data (81.4\%)
outperforms the E2E model at 100\% data (75.4\%) by nearly 6 points. Foundation model
features provide a higher accuracy floor at every data fraction. In low-data settings
like VizWiz, CLIP embeddings are a far more data-efficient starting point.

% ════════════════════════════════════════════════════════════
\clearpage
\appendix
\section*{Appendix: Code}
% ════════════════════════════════════════════════════════════

\subsection*{A.1 End-to-End Models (Challenges 1 \& 2): \texttt{run\_e2e.py}}

\lstinputlisting[language=Python]{run_e2e.py}

\subsection*{A.2 CLIP Feature-Based Models (Challenges 3 \& 4): \texttt{solution\_clip.py}}

\lstinputlisting[language=Python]{solution_clip.py}

\subsection*{A.3 Controlled Experiments (Sections 1 \& 2 CLIP): \texttt{run\_experiments.py}}

\lstinputlisting[language=Python]{run_experiments.py}

\subsection*{A.4 Data Efficiency Experiments (Section 3): \texttt{data\_efficiency.py}}

\lstinputlisting[language=Python]{data_efficiency.py}

\end{document}
